\documentclass[11pt]{article}

\usepackage[a4paper,margin=1in]{geometry}
\usepackage{hyperref}
\usepackage{listings}
\usepackage{xcolor}
\usepackage{enumitem}
\usepackage{titlesec}

\titleformat{\section}{\large\bfseries}{\thesection}{1em}{}
\titleformat{\subsection}{\normalsize\bfseries}{\thesubsection}{1em}{}

\definecolor{codegray}{rgb}{0.95,0.95,0.95}
\definecolor{commentgreen}{rgb}{0.0,0.5,0.0}

\lstset{
  basicstyle=\ttfamily\small,
  backgroundcolor=\color{codegray},
  frame=single,
  breaklines=true,
  showstringspaces=false,
  keywordstyle=\color{blue},
  commentstyle=\color{commentgreen},
  stringstyle=\color{red}
}

\title{\textbf{Lab Standard: SLURM Scripts for Scalable ML Experiments}}
\author{Medical Computer Vision Lab (MCV Lab)}
\date{\today}

\begin{document}
\maketitle

\section{Purpose}
This document defines the lab standard for running ML experiments on HPC clusters using SLURM.
All jobs must be:
\begin{itemize}[nosep]
    \item reproducible
    \item logged
    \item checkpointed
    \item sweep-ready
\end{itemize}

\section{General Rules}
\begin{enumerate}[label=\textbf{S\arabic*.}, leftmargin=*]
    \item \textbf{Never edit training scripts in HPC jobs}. Use configs/overrides only.
    \item \textbf{Always log SLURM environment variables}.
    \item \textbf{Always write outputs into a run folder}.
    \item \textbf{Always checkpoint (last + best)}.
    \item \textbf{Always specify GPU, CPU, RAM explicitly}.
\end{enumerate}

\section{Folder Standard}
SLURM logs must be stored separately from run outputs:

\begin{lstlisting}
project/
  runs/
  slurm_logs/
    %x_%j.out
    %x_%j.err
  scripts/
    run_train.slurm
\end{lstlisting}

\section{Baseline SLURM Job (Single Run)}
Save as \texttt{scripts/run\_train.slurm}:

\begin{lstlisting}[language=bash]
#!/bin/bash
#SBATCH --job-name=mcv_train
#SBATCH --output=slurm_logs/%x_%j.out
#SBATCH --error=slurm_logs/%x_%j.err
#SBATCH --time=08:00:00
#SBATCH --partition=gpu
#SBATCH --gres=gpu:1
#SBATCH --cpus-per-task=8
#SBATCH --mem=32G

# Optional: email notifications
# SBATCH --mail-type=END,FAIL
# SBATCH --mail-user=your_email@uvic.ca

set -euo pipefail

echo "=== SLURM JOB INFO ==="
echo "JOB_ID:      ${SLURM_JOB_ID}"
echo "JOB_NAME:    ${SLURM_JOB_NAME}"
echo "NODELIST:    ${SLURM_NODELIST}"
echo "GPUS:        ${SLURM_GPUS:-unknown}"
echo "CPUS:        ${SLURM_CPUS_PER_TASK}"
echo "WORKDIR:     ${SLURM_SUBMIT_DIR}"
echo "START_TIME:  $(date)"

cd ${SLURM_SUBMIT_DIR}

# Activate environment
# Option A: conda
# source ~/miniconda3/etc/profile.d/conda.sh
# conda activate mcv

# Option B: module-based
# module load python/3.10 cuda/12.1
# pip install -r requirements.txt

nvidia-smi || true

# Run training
python train.py \
  --dataset rsna_trauma \
  --model resnet50 \
  --lr 1e-4 \
  --seed 0 \
  --amp \
  --out_dir runs \
  --run_name "rsna_resnet50_seed0_lr1e4_${SLURM_JOB_ID}"

echo "END_TIME: $(date)"
\end{lstlisting}

\section{SLURM Array for Sweeps}
SLURM arrays are the lab standard for running sweeps.

\subsection{Example: Sweep over seeds}
Save as \texttt{scripts/sweep\_seeds.slurm}:

\begin{lstlisting}[language=bash]
#!/bin/bash
#SBATCH --job-name=mcv_seed_sweep
#SBATCH --output=slurm_logs/%x_%A_%a.out
#SBATCH --error=slurm_logs/%x_%A_%a.err
#SBATCH --time=08:00:00
#SBATCH --partition=gpu
#SBATCH --gres=gpu:1
#SBATCH --cpus-per-task=8
#SBATCH --mem=32G
#SBATCH --array=0-4

set -euo pipefail
cd ${SLURM_SUBMIT_DIR}

SEEDS=(0 1 2 3 4)
SEED=${SEEDS[$SLURM_ARRAY_TASK_ID]}

python train.py \
  --dataset rsna_trauma \
  --model resnet50 \
  --lr 1e-4 \
  --seed ${SEED} \
  --amp \
  --out_dir runs \
  --run_name "rsna_resnet50_seed${SEED}_job${SLURM_JOB_ID}"
\end{lstlisting}

Submit:
\begin{lstlisting}
sbatch scripts/sweep_seeds.slurm
\end{lstlisting}

\section{Hydra + SLURM Recommended Pattern}
If using Hydra multirun, we prefer either:
\begin{itemize}[nosep]
    \item Hydra \texttt{-m} multirun inside a single SLURM job (small sweeps)
    \item SLURM array + Hydra overrides (large sweeps; recommended)
\end{itemize}

\subsection{SLURM array + Hydra override template}
\begin{lstlisting}[language=bash]
python train.py \
  dataset=rsna_trauma \
  model=resnet50 \
  optim=adamw \
  optim.lr=1e-4 \
  seed=${SEED}
\end{lstlisting}

\section{Multi-GPU Training (DDP)}
\subsection{Standard: torchrun}
Save as \texttt{scripts/ddp\_train.slurm}:

\begin{lstlisting}[language=bash]
#!/bin/bash
#SBATCH --job-name=mcv_ddp
#SBATCH --output=slurm_logs/%x_%j.out
#SBATCH --error=slurm_logs/%x_%j.err
#SBATCH --time=12:00:00
#SBATCH --partition=gpu
#SBATCH --gres=gpu:2
#SBATCH --cpus-per-task=12
#SBATCH --mem=64G

set -euo pipefail
cd ${SLURM_SUBMIT_DIR}

export OMP_NUM_THREADS=${SLURM_CPUS_PER_TASK}

# Use torchrun for DDP
torchrun --nproc_per_node=2 train.py \
  --dataset rsna_trauma \
  --model resnet50 \
  --lr 1e-4 \
  --seed 0 \
  --amp \
  --out_dir runs \
  --run_name "rsna_ddp2_${SLURM_JOB_ID}"
\end{lstlisting}

\section{Checklist for SLURM Job Reviews}
Before using a SLURM script in the lab:
\begin{itemize}
    \item Logs go to \texttt{slurm\_logs/}
    \item Run outputs go to \texttt{runs/}
    \item Script logs SLURM metadata
    \item Script runs without editing python source
    \item Seeds are parameterized
    \item Checkpoints saved (best + last)
\end{itemize}

\end{document}
